% !TEX root = main.tex
\section{Misc Notations}\label{sec:misc-notations}
Sometimes we will construct a new set from two given set by some operations. Here we define two kind of constructions.
\begin{definition}
    We denote $A + b = \{ x + b \mid x \in A \}$ and $aA = \{ ax \mid x \in A \}$,
    denote $A \oplus B = \{ x + y \mid x \in A, y \in B \}$ and $A \otimes B = \{ xy \mid x \in A, y \in B \}$.
    Finally, we denote $\{ A < s \} = \{ x \in A \mid x < s \}$; the definitions for the cases involving $>, \le,$ and $\ge$ are analogous.
\end{definition}
\begin{proposition}
    $A \oplus B = B \oplus A$ and $A \otimes B = B \otimes A$.
\end{proposition}
If $\forall s\in C = A \otimes B$, there exists a unique pair $x \in A$ and $y \in B$ such that $s = xy$, then $C$ is said to multiplily decompose into $A$ and $B$ uniqely. If $\forall s\in D = A \otimes B$, there exists a unique pair $x \in A$ and $y \in B$ such that $s = x + y$, then $D$ is said to additively decompose into $A$ and $B$ uniqely.

\begin{definition}
    Suppose $R$ be a ring, then we denote $R^{*} = R - \{0\}$.
\end{definition}
\begin{definition}
Given $G$ be a group and $H\le G$.
We denote $G/H$ or $\frac{G}{H}$ as the quotioniant group of $G$ relative to $H$.
\end{definition}
